% Options for packages loaded elsewhere
\PassOptionsToPackage{unicode}{hyperref}
\PassOptionsToPackage{hyphens}{url}
\documentclass[
]{article}
\usepackage{xcolor}
\usepackage{amsmath,amssymb}
\setcounter{secnumdepth}{-\maxdimen} % remove section numbering
\usepackage{iftex}
\ifPDFTeX
  \usepackage[T1]{fontenc}
  \usepackage[utf8]{inputenc}
  \usepackage{textcomp} % provide euro and other symbols
\else % if luatex or xetex
  \usepackage{unicode-math} % this also loads fontspec
  \defaultfontfeatures{Scale=MatchLowercase}
  \defaultfontfeatures[\rmfamily]{Ligatures=TeX,Scale=1}
\fi
\usepackage{lmodern}
\ifPDFTeX\else
  % xetex/luatex font selection
\fi
% Use upquote if available, for straight quotes in verbatim environments
\IfFileExists{upquote.sty}{\usepackage{upquote}}{}
\IfFileExists{microtype.sty}{% use microtype if available
  \usepackage[]{microtype}
  \UseMicrotypeSet[protrusion]{basicmath} % disable protrusion for tt fonts
}{}
\makeatletter
\@ifundefined{KOMAClassName}{% if non-KOMA class
  \IfFileExists{parskip.sty}{%
    \usepackage{parskip}
  }{% else
    \setlength{\parindent}{0pt}
    \setlength{\parskip}{6pt plus 2pt minus 1pt}}
}{% if KOMA class
  \KOMAoptions{parskip=half}}
\makeatother
\setlength{\emergencystretch}{3em} % prevent overfull lines
\providecommand{\tightlist}{%
  \setlength{\itemsep}{0pt}\setlength{\parskip}{0pt}}
\usepackage{bookmark}
\IfFileExists{xurl.sty}{\usepackage{xurl}}{} % add URL line breaks if available
\urlstyle{same}
\hypersetup{
  hidelinks,
  pdfcreator={LaTeX via pandoc}}

\author{}
\date{}

\begin{document}

\section{Inference from an L1 and L2
Perspective}\label{inference-from-an-l1-and-l2-perspective}

\textbf{Author:} Jeremy H. Carroll \textbf{Date:} March 2026
\textbf{Version:} v1.0 (Pre-Print)

\subsection{Abstract}\label{abstract}

In the foundational \(\ell^1\) verification framework, inference is
fundamentally redefined. It is not solely the probabilistic prediction
of the next token, as seen in standard \(\ell^2\) large language models.
Rather, it can be modeled as a \textbf{morphism} in a category of
reasoning states. This paper defines the precise topological mechanics
of autonomous inference, contrasting the continuous \(\ell^2\)
generation of token probabilities against the discrete \(\ell^1\)
embedding of causal reasoning steps into a globally commutating directed
acyclic graph.

\begin{center}\rule{0.5\linewidth}{0.5pt}\end{center}

\subsection{\texorpdfstring{1. Inference as a Morphism
(\texttt{InferenceStep})}{1. Inference as a Morphism (InferenceStep)}}\label{inference-as-a-morphism-inferencestep}

An individual step of inference acts as a directed arrow that transforms
one \texttt{ReasoningState} (representing the source context, domain,
and constraint boundary) into a new \texttt{ReasoningState} (the
target).

Every inference step is exclusively executed by a specific Lambda
operator and carries a defined \texttt{MorphismKind} (such as a
computation, generation, or verification step). By modeling reasoning
categorially, every operational move an AI takes is mathematically
typed, tracked, and structurally composable with subsequent or preceding
steps.

\begin{center}\rule{0.5\linewidth}{0.5pt}\end{center}

\subsection{2. Topological Validation via
Commutativity}\label{topological-validation-via-commutativity}

Because inference is treated as a sequence of formal morphisms, its
validity is tested geometrically using \textbf{commutative diagrams}.

If the system applies two different chains of reasoning to navigate from
concept A to concept C---for example, taking a direct inductive path
\(h\) versus a two-step deductive path \(g \circ f\)---the results must
perfectly match: \[(g \circ f)(a) \equiv h(a)\]

If the diagram fails to commute, it acts as a topological tear detector,
proving that two distinct parts of the system mathematically disagree
about objective reality. In an \(\ell^2\) paradigm, non-commuting logic
is statistically averaged out. In the \(\ell^1\) paradigm, the failure
to commute immediately spikes the coboundary defect, halting the
assumption that the logic is coherent.

\subsubsection{\texorpdfstring{2.1 Theorem (Non-Commutativity Implies
\(\ell^1\) Defect Lower
Bound)}{2.1 Theorem (Non-Commutativity Implies \textbackslash ell\^{}1 Defect Lower Bound)}}\label{theorem-non-commutativity-implies-ell1-defect-lower-bound}

Let \(\mathcal{P}\) be a causal poset, \(s \in C^0(\mathcal{P}, F)\) a
section, \(\delta = d_0 s\) the \(\ell^1\) coboundary, and
\(I(s) = \|\delta\|_1\). Consider a commutative diagram candidate
\(A \xrightarrow{f} B \xrightarrow{g} C\) and \(A \xrightarrow{h} C\).
Define the \textbf{commutativity defect}:
\[\Delta(a) := (g \circ f)(a) - h(a)\]

\textbf{Claim:} If the diagram does not commute (\(\Delta(a) \neq 0\)),
then \(I(s) \geq \|\Delta(a)\|_1\).

\textbf{Proof:} Each morphism corresponds to an edge path in the poset
(\(p_1: A \to B \to C\) and \(p_2: A \to C\)). The coboundary
accumulates along paths such that
\(\sum_{e \in p_1} \delta_e(s) = (g \circ f)(a)\) and
\(\sum_{e \in p_2} \delta_e(s) = h(a)\). Subtracting these yields a
cycle:
\[\sum_{e \in p_1} \delta_e(s) - \sum_{e \in p_2} \delta_e(s) = \Delta(a)\]
By the triangle inequality:
\[\|\Delta(a)\|_1 \leq \sum_{e \in p_1 \cup p_2} \|\delta_e(s)\|_1 \leq I(s)\]

\textbf{Conclusion:} Non-commutativity forces a strictly positive
\(\ell^1\) defect lower bound. Commutativity equates to zero defect
locally; non-commutativity ensures a structurally irreducible defect
cost.

\begin{center}\rule{0.5\linewidth}{0.5pt}\end{center}

\subsection{3. The Exact Mathematical Definition of a
Hallucination}\label{the-exact-mathematical-definition-of-a-hallucination}

When an inference step produces an unresolvable contradiction (a
non-zero \(\ell^1\) coboundary defect), it fails to structurally embed
into the global constraint structure of the domain graph.

In this categorical framework, \textbf{a hallucination is formally
defined as an inference step that cannot validly embed into the causal
poset}. It is not a statistical anomaly or ``bad data''; it is a massive
geometric breakdown. The generated edge physically violates the
commutative geometry of the object-level reasoning space.

\begin{center}\rule{0.5\linewidth}{0.5pt}\end{center}

\subsection{\texorpdfstring{4. Branching and the
\texttt{ReasoningGraph}}{4. Branching and the ReasoningGraph}}\label{branching-and-the-reasoninggraph}

As inference progresses, it generates a \texttt{ReasoningGraph}---a
directed acyclic graph (DAG) where nodes are reasoning artifacts
(prompts, outputs, refinements) and edges encode the formal parent-child
derivation.

When the system faces uncertainty and diverges to explore multiple
inference paths simultaneously, this branching transforms the standard
reasoning DAG into \textbf{a stratified branching structure with
cohomological obstruction depth}, suggestive of derived stack
constructions. In this mathematically layered space, the cohomological
depth (the \(H^1\) measurement) explicitly tracks how many of these
branches are structurally obstructed from gluing back together into a
single, globally consistent truth.

\begin{center}\rule{0.5\linewidth}{0.5pt}\end{center}

\subsection{5. Resolution of Obstructed
Inference}\label{resolution-of-obstructed-inference}

If an inference path hits a dead end (measuring an \(H^1 > 0\) geometric
obstruction), the \(\ell^1\) system does not average out the error to
find a ``phantom point.'' Instead, it natively resolves the conflict by
relying on an explicit mechanism:

\subsubsection{5.1 Autopoietic Control}\label{autopoietic-control}

If an objective obstruction is identified, the system triggers the
Autopoietic Iterator. This metacognitive control mechanism inserts new
structural constraints (a specific \texttt{ReasoningMove}) designed to
bridge the unresolvable gap. This operation corresponds to graph
augmentation steps that iteratively reduce the defect until systemic
consistency is achieved.

\subsubsection{5.2 Theorem (Autopoietic Resolution of
Non-Commutativity)}\label{theorem-autopoietic-resolution-of-non-commutativity}

Every non-commuting diagram in the modeled reasoning graph induces a
positive \(\ell^1\) defect. Under Autopoietic dynamics, the iterator is
mathematically forced to act on non-commuting diagrams to reduce the
defect: - \textbf{Repairable (Case A):} There exists an expansion
\(\Lambda\) such that
\(I(\mathcal{P} \cup \Lambda, s) < I(\mathcal{P}, s)\). The iterator
inserts edges, forcing the diagram to become commutative. -
\textbf{Topologically Protected (Case B):} No such \(\Lambda\) exists
(\(I^* = \Phi([\delta]) > 0\)); non-commutativity persists and becomes a
stable structural invariant.

\subsubsection{5.3 Structural Expansion vs.~Fixed-Architecture
Optimization}\label{structural-expansion-vs.-fixed-architecture-optimization}

Unlike fixed-architecture optimization systems (Standard ML, MCTS, RL)
that optimize values or policies over a fixed state space and handle
contradictions statistically, the Autopoietic Iterator performs descent
over both state and structure. It resolves inconsistency by modifying
the underlying causal graph (pushout) rather than averaging over it.

\begin{center}\rule{0.5\linewidth}{0.5pt}\end{center}

\subsection{6. Monotone Dynamics of
Inference}\label{monotone-dynamics-of-inference}

In this topological framework, the \(\ell^1\) coboundary norm defines a
monotone quantity analogous to entropy gradients in physical systems.
The \(\ell^1\) coboundary norm formally measures topological tension
across the \texttt{Reasoning\ Graph}, ensuring inference is
mathematically constrained by monotonic descent principles.

The Omega Engine executes this explicitly across four internal modules:

\subsubsection{6.1 Epistemic Tension as Sheaf
Restrictions}\label{epistemic-tension-as-sheaf-restrictions}

Abstract inference states are modeled as local sections of a pre-sheaf
over the inference phase space. The \(\ell^1\) coboundary norm
calculates the absolute mismatch across each overlapping edge in the
complex. * \textbf{The Calculation:} \(I(s) = \sum |\delta_e(s)|\). *
\textbf{Interpretation:} This \(\ell^1\) tension measures structural
disequilibrium---the exact degree to which local claims fail to glue
into a globally coherent macrostate.

\subsubsection{\texorpdfstring{6.2 The Cohomological Obstruction
(\texttt{stat\_mech.rs})}{6.2 The Cohomological Obstruction (stat\_mech.rs)}}\label{the-cohomological-obstruction-stat_mech.rs}

When the system projects a highly complex branching state into a single
verifiable inference path, it evaluates the exact information
differential. * \textbf{The Calculation:} \(H^1\) obstruction scaling. *
\textbf{Interpretation:} This evaluates the ``information cost'' of a
reasoning projection. The \(\ell^1\) obstruction here is mathematically
guaranteed to be \(\ge 0\), representing the irreducible tension that
prevents the AI from collapsing contradictory logics into a localized
point.

\subsubsection{\texorpdfstring{6.3 Inference Trace Evaluation
(\texttt{compute\_cognitive\_z.rs})}{6.3 Inference Trace Evaluation (compute\_cognitive\_z.rs)}}\label{inference-trace-evaluation-compute_cognitive_z.rs}

The engine can evaluate individual epistemic traces by treating them as
weighted paths in the reasoning DAG. The ``action'' equivalent \(S[s]\)
of a specific inference sequence is defined rigorously as its total
\(\ell^1\) topological defect: \(S[s] = \|\delta s\|_1\). The logged
\texttt{perception.entropy} metric serves as a heuristic evaluation
proxy for tracking this coboundary cost during generation.

\subsubsection{\texorpdfstring{6.4 Coupled within the Lyapunov Control
Function
(\texttt{crf\_tensor.rs})}{6.4 Coupled within the Lyapunov Control Function (crf\_tensor.rs)}}\label{coupled-within-the-lyapunov-control-function-crf_tensor.rs}

The system bounds active generation using a Lyapunov-analogous descent
function that restricts generation to monotonic pathways. * \textbf{The
Calculation:} \(V(s) \propto I(s) + \gamma (\text{divergence bounds})\).
* \textbf{Interpretation:} The control function bounds the maximum
random topological diffusion permitted by the system. If divergence
spirals (\(V > 0\)), topological negative feedback forces compression
along the constraint graph.

Hallucination in structural intelligence is a \textbf{cohomological
obstruction}. The inference steps provide the 0-cochain values (the
nodes), and the \(\ell^1\) norm mathematically measures the consistency
limits across the causal edges.

\begin{center}\rule{0.5\linewidth}{0.5pt}\end{center}

\subsection{\texorpdfstring{7. Calculating Inference: \(\ell^1\)
vs.~\(\ell^2\)}{7. Calculating Inference: \textbackslash ell\^{}1 vs.~\textbackslash ell\^{}2}}\label{calculating-inference-ell1-vs.-ell2}

The core difference between calculating inference in \(\ell^1\) geometry
versus standard \(\ell^2\) geometry comes down to \textbf{sparsity
vs.~diffusion}.

\begin{itemize}
\tightlist
\item
  \textbf{\(\ell^1\) Inference (Sparsity \& Isolation):} In \(\ell^1\)
  geometry, inference errors are evaluated by computing the exact
  coboundary defect \(I(s) = \sum \|\delta_e(s)\|_1\), which sums the
  absolute discrepancies across all reasoning edges. When resolving
  errors, \(\ell^1\) calculation finds the \textbf{median} rather than
  the mean, ensuring that the system does not interpolate between
  contradictory claims. It encourages the sparsity of data and forces
  unresolved contradictions to concentrate onto lower-dimensional
  boundaries (min-cuts), structurally localizing inconsistency.
\item
  \textbf{\(\ell^2\) Inference (Diffusion \& Smoothing):} \(\ell^2\)
  inference evaluates errors using the sum of squares, inherently
  allowing contradictions of opposite signs to partially cancel each
  other out. In the Rust engine, \(\ell^2\) optimization behaves like a
  graph Laplacian or a heat equation; it diffuses the residual error
  globally across the entire network, smoothing out gradients but
  mathematically permitting the diffusion of inconsistency across
  structural boundaries.
\item
  \textbf{Calculating the Difference (The Tension):} The mathematical
  difference between the two is explicitly calculated using the
  \textbf{Inflation Ratio}. The system evaluates the raw defect using an
  \(\ell^2\) Hodge decomposition (splitting it into exact and harmonic
  components via orthogonal projection) and compares the inflated cost
  against the true \(\ell^1\) quotient norm \(\Phi\). Mapping a sparse
  exact defect into an \(\ell^2\) harmonic space universally inflates
  the mathematical cost, strictly approaching a ratio of 3 as observed
  in cyclic graph constructions (see Paper 007b).
\end{itemize}

\begin{center}\rule{0.5\linewidth}{0.5pt}\end{center}

\subsection{8. Mathematical Tools in the Rust
Codebase}\label{mathematical-tools-in-the-rust-codebase}

To perform these evaluations efficiently, the \texttt{omega\_core}
codebase utilizes several rigorous mathematical and topological engines:

\subsubsection{\texorpdfstring{8.1 The \texttt{CoboundaryNorm} Engine
(\texttt{l028})}{8.1 The CoboundaryNorm Engine (l028)}}\label{the-coboundarynorm-engine-l028}

This is the foundational tool for \(\ell^1\) inference. It takes a
\texttt{FinitePoset} and a \texttt{Section} (the values assigned to the
nodes) and strictly computes \texttt{compute\_l1\_rigidity()}, natively
generating the total absolute defect across all covering edges.

\subsubsection{\texorpdfstring{8.2 \(\ell^1\)-Native Hodge Decomposition
\& ISTA
(\texttt{l302})}{8.2 \textbackslash ell\^{}1-Native Hodge Decomposition \& ISTA (l302)}}\label{ell1-native-hodge-decomposition-ista-l302}

Instead of relying on the classical \(\ell^2\) Hodge projection---which
uses Euclidean inner products---the codebase implements
\texttt{L1NativeDecomposition}. To compute the primal minimizer
(\(\min_f \|\delta - d_0 f\|_1\)), it uses \texttt{solve\_primal\_ista},
an Iterative Soft-Thresholding Algorithm that natively handles
\(\ell^1\) projection to gracefully isolate exact gauge artifacts from
irreducible topological tears.

\subsubsection{\texorpdfstring{8.3 The Autopoietic Iterator /
Wasserstein Descent
(\texttt{l166})}{8.3 The Autopoietic Iterator / Wasserstein Descent (l166)}}\label{the-autopoietic-iterator-wasserstein-descent-l166}

To calculate the structural evolution of inference, this tool implements
discrete Wasserstein-1 gradient descent. Instead of global \(\ell^2\)
smoothing, it iteratively detects the worst \(\ell^1\) tear, calculates
local \(\ell^1\)-optimal medians on star neighborhoods, and surgically
augments the poset with new geometric edges.

\subsubsection{\texorpdfstring{8.4 Linear Programming (LP) and Network
Simplex
(\texttt{l097})}{8.4 Linear Programming (LP) and Network Simplex (l097)}}\label{linear-programming-lp-and-network-simplex-l097}

To definitively calculate the irreducible domain complexity floor
(\(\Phi\)), the codebase frames the \(\ell^1\) quotient norm as a Linear
Program (LP). The LP solver constructs the exact primal-dual pairing by
evaluating bounded divergence-free circulations against the defect
field, calculating the ultimate topological obstruction in polynomial
time \(O(|E|^{1.5})\).

\subsubsection{\texorpdfstring{8.5 Spectral Graph Laplacian
(\texttt{l097})}{8.5 Spectral Graph Laplacian (l097)}}\label{spectral-graph-laplacian-l097}

When specifically measuring \(\ell^2\) geometries (or evaluating
precisely why they structurally fail), the system uses the
\texttt{GraphLaplacian} (\(L = D - A\)) to calculate the Dirichlet
energy \(E(\rho) = Tr(\rho L)\). It applies \texttt{spectral\_diffusion}
(the heat kernel \(e^{-tL}\)) to simulate how \(\ell^2\) mechanisms
smoothly spread probability states across graph neighborhoods. It also
utilizes the \texttt{lanczos} algorithm to efficiently extract top-k
eigenvalues for extremely fast von Neumann entropy approximations.

\begin{center}\rule{0.5\linewidth}{0.5pt}\end{center}

\subsection{9. Conclusion}\label{conclusion}

Inference in the verification framework is the continuous, dynamic
mapping of morphisms across a causal poset, tightly governed by the
\(\ell^1\) coboundary norm to ensure every logical step mathematically
aligns with objective thermodynamic reality. This demonstrates
structural limitations of \(\ell^2\) autoregressive inference in the
presence of contradictory constraints, motivating \(\ell^1\)-based
alternatives.

\end{document}
